\subsection{Milestone 2: Model Building -- Reproducibility}
\label{sec:m2}
\guidance{
  \textbf{Template:} Describe how the different software engineering practices and tools have been applied.\\
  EDN entries here only for reproducibility decisions that go beyond what the course prescribes.
}

\subsubsection{Project structure}
\label{sec:m2-structure}
\guidance{
  \textbf{Template:} e.g.\ Cookiecutter.\\
  \textbf{Rubric (5 pts):} How well does the project structure support clarity, maintainability, and adaptation to the project's needs?\\
  We use the Cookiecutter Data Science template \cite{ccds} with uv as dependency manager.
  Justify any deviations from it (the course allows justified deviations).
}

We use the Cookiecutter Data Science template \cite{ccds} with \texttt{uv} as dependency manager, and deviate from it in three places.
The course allows justified deviations; each of ours is recorded in the EDN.

\paragraph{One module per pipeline stage (EDN-28).}
The template generates a flat \texttt{dataset.py}, \texttt{features.py} and \texttt{modeling/train.py}.
We replaced those with one module per DVC stage, under \texttt{recommenditos/data/} and \texttt{recommenditos/modeling/}, which is also the layout the teachers' demo repository uses.
The technical argument is the dependency graph: a stage's \texttt{deps} should name the code that stage actually runs, so \texttt{dvc repro} reruns the minimum.
A flat layout cannot express that, because one file backs several stages.
The organisational argument is that the pipeline is sequential, so parallel work is only possible if each ticket owns its own files.

\paragraph{A \texttt{params.yaml} instead of constants in \texttt{config.py} (EDN-30).}
The demo keeps the seed, the split ratios and the hyperparameters as Python constants.
We moved everything configurable into \texttt{params.yaml} and left only paths in \texttt{config.py}.
A parameter change then reruns exactly the stages that read it, \texttt{dvc params diff} shows what changed between two commits, and \texttt{dvc exp} becomes usable.
Paths stay in code so that nothing in the project depends on a machine-specific absolute path, which is a wart the demo's \texttt{config.py} does have.

\paragraph{Two stages the demo does not have.}
\texttt{features} exists because the demo's text model needs no feature engineering and our tabular model does.
\texttt{evaluate} exists because the demo asserts its metric threshold inside \texttt{tests/test\_model.py}, whereas NFR-01 gates deployment on six success criteria and EDN-12 makes promotion a human decision, so someone has to be able to read the result off an artefact.

We also added \texttt{recommenditos/schema.py}, the structural contract for the processed data, which the stages code against instead of against each other (EDN-31).

\subsubsection{Code versioning}
\label{sec:m2-code-versioning}
\guidance{
  \textbf{Template:} e.g.\ GitHub Flow.\\
  \textbf{Rubric (15 pts, shared with data versioning):} How do our version control practices keep collaboration smooth and minimise mistakes?\\
  Already in place: GitHub Flow with \texttt{feature/} and \texttt{fix/} branches, branch protection on \texttt{main}, squash merge as the only merge method, Conventional Commits enforced by pre-commit and by a CI check on PR titles, pre-commit hooks (ruff, nbstripout, file hygiene).
}

\tbd{code versioning workflow}

\subsubsection{Data versioning}
\label{sec:m2-data-versioning}
\guidance{
  \textbf{Template:} e.g.\ DVC.\\
  \textbf{Rubric:} How effectively does our data versioning strategy guarantee reproducibility and scalability?\\
  The demo uses DVC with a DagsHub remote (HTTP option) and a \texttt{dvc.yaml} pipeline \texttt{download -> preprocess -> validate -> split -> train} with code files as dependencies; commit \texttt{dvc.lock}.
  Consider a pipeline figure.
}

\tbd{data versioning narrative and the pipeline figure (\texttt{dvc dag -{}-md})}

The pipeline is \texttt{download \textrightarrow{} preprocess \textrightarrow{} \{configure\_gx, validate-data\} \textrightarrow{} split \textrightarrow{} features \textrightarrow{} train \textrightarrow{} evaluate}, with \texttt{features} and \texttt{train} iterating a mapping in \texttt{params.yaml} through \texttt{foreach}, so their stages are named after the item (\texttt{train@lgbm-basic}) rather than its position.
Unlike the demo's, our \texttt{validate-data} stage fails the pipeline when an expectation fails instead of only logging the failure count.

\subsubsection{Experiment tracking}
\label{sec:m2-experiment-tracking}
\guidance{
  \textbf{Template:} e.g.\ MLflow.\\
  \textbf{Rubric (5 pts):} How does our experiment tracking help us and others compare, understand, and extend our results?\\
  The demo uses DagsHub as shared MLflow tracking server with credentials in \texttt{.env}.
  Show what we log (parameters, metrics, artifacts, emissions) and a comparison of runs.
}

\tbd{experiment tracking setup and results}
